\documentclass[10pt,a5paper,landscape]{article}
\usepackage[french]{babel}
\usepackage{fontspec}
\usepackage{geometry}
\usepackage{graphicx}
\usepackage{xcolor}
\usepackage{fancyhdr}
\usepackage{hyperref}
\usepackage{microtype}
\usepackage{ragged2e}
\usepackage{multicol}
\usepackage{tikz}
\usetikzlibrary{calc}

\geometry{left=1cm,right=1cm,top=0.85cm,bottom=0.8cm,footskip=0.35cm}
\setlength{\parindent}{0pt}
\setlength{\parskip}{0.35em}
\setlength{\emergencystretch}{2em}
\setlength{\columnsep}{0.9cm}
\hyphenpenalty=8500
\exhyphenpenalty=8500

\definecolor{ink}{HTML}{111111}
\definecolor{paper}{HTML}{F3F1EC}
\definecolor{acid}{HTML}{D7FF37}
\definecolor{muted}{HTML}{777570}
\definecolor{soft}{HTML}{B6B2AA}
\definecolor{rulegray}{HTML}{D8D4CB}
\pagecolor{paper}
\color{ink}

\setmainfont{Satoshi-Regular.otf}[Path=../font/,BoldFont=Satoshi-Bold.otf,ItalicFont=Satoshi-Italic.otf]
\newfontfamily\displayfont{Satoshi-Black.otf}[Path=../font/]
\newfontfamily\mediumfont{Satoshi-Medium.otf}[Path=../font/,BoldFont=Satoshi-Bold.otf]
\newfontfamily\lightfont{Satoshi-Light.otf}[Path=../font/,BoldFont=Satoshi-Bold.otf,ItalicFont=Satoshi-Italic.otf]

\hypersetup{colorlinks=true,linkcolor=ink,urlcolor=ink,pdfauthor={Alexandre Mathieu},pdftitle={Portfolio V2 2018-2026}}
\pagestyle{fancy}
\fancyhf{}
\renewcommand{\headrulewidth}{0pt}
\renewcommand{\footrulewidth}{0pt}
\fancyfoot[L]{\fontsize{5.3}{6.4}\selectfont\lightfont\color{muted} ALEXANDRE MATHIEU / PORTFOLIO}
\fancyfoot[R]{\fontsize{5.3}{6.4}\selectfont\lightfont\color{muted} \thepage}

\newcommand{\hairline}{\textcolor{rulegray}{\rule{\linewidth}{0.25pt}}}
\newcommand{\fullBleedImage}[1]{%
  \begin{tikzpicture}[remember picture,overlay]
    \node[anchor=center,inner sep=0] at (current page.center)
      {\includegraphics[width=\paperwidth,height=\paperheight]{#1}};
  \end{tikzpicture}%
}

\newcommand{\projectMetaItem}[2]{%
  \begin{minipage}[t]{0.47\linewidth}
    {\fontsize{4.9}{5.8}\selectfont\mediumfont\color{muted}\MakeUppercase{#1}}\par
    \vspace{0.04cm}
    {\fontsize{6.2}{7.3}\selectfont\lightfont\color{ink}#2}\par
    \vspace{0.20cm}
  \end{minipage}\hfill
}
\newcommand{\leadText}[1]{{\fontsize{8.2}{9.8}\selectfont\mediumfont\color{ink}#1}\par\vspace{0.25cm}}
\newcommand{\projectSubheading}[1]{\vspace{0.11cm}{\fontsize{6.3}{7.6}\selectfont\mediumfont\color{ink}\MakeUppercase{#1}}\par\vspace{-0.02cm}}

\newcommand{\portfolioTocEntry}[3]{%
  \hyperlink{project-#1}{%
    \noindent
    \begin{minipage}[c]{0.10\linewidth}{\fontsize{6.2}{7}\selectfont\mediumfont\color{white!50}#1}\end{minipage}
    \begin{minipage}[c]{0.70\linewidth}{\fontsize{8.8}{9.7}\selectfont\mediumfont\color{white}#2}\end{minipage}
    \begin{minipage}[c]{0.16\linewidth}\raggedleft{\fontsize{5.8}{6.8}\selectfont\lightfont\color{white!50}#3}\end{minipage}
  }\par\vspace{0.10cm}
  \textcolor{white!18}{\rule{\linewidth}{0.2pt}}\par\vspace{0.10cm}
}

\newcommand{\projectIntroPage}[5]{%
  \clearpage\thispagestyle{empty}\hypertarget{project-#1}{}
  \ifx&#5&\else\fullBleedImage{#5}\fi
  \begin{tikzpicture}[remember picture,overlay]
    \fill[ink,opacity=0.15] (current page.south west) rectangle (current page.north east);
    \fill[ink,opacity=0.92] (current page.south west) rectangle ([xshift=12.2cm,yshift=4.15cm]current page.south west);
    \node[anchor=north west,text=white,inner sep=0,text width=10.7cm]
      at ([xshift=0.8cm,yshift=3.55cm]current page.south west) {%
        {\fontsize{5.5}{6.5}\selectfont\mediumfont\color{acid}PROJET #1 \hfill #3}\par
        \vspace{0.28cm}
        {\fontsize{21}{20.5}\selectfont\displayfont\color{white}#2}\par
        \vspace{0.20cm}
        {\fontsize{6.2}{7.3}\selectfont\lightfont\color{white!72}#4}
      };
    \node[anchor=south east,rotate=90,text=white,inner sep=0]
      at ([xshift=-0.33cm,yshift=0.55cm]current page.south east)
      {\fontsize{4.8}{5.6}\selectfont\mediumfont IMAGE / TERRITOIRE / MATIERE};
  \end{tikzpicture}\null
}

\long\def\projectStoryPage#1#2#3#4#5{%
  \clearpage\thispagestyle{empty}
  \begin{tikzpicture}[remember picture,overlay]
    \fill[paper] (current page.south west) rectangle (current page.north east);
    \ifx&#5&\else
      \node[anchor=north east,inner sep=0] at (current page.north east)
        {\includegraphics[width=9.7cm,height=14.8cm]{#5}};
    \fi
    \fill[acid] ([xshift=-9.7cm]current page.north east) rectangle ([xshift=-9.52cm]current page.south east);
    \node[anchor=north west,inner sep=0,text width=9.45cm]
      at ([xshift=0.75cm,yshift=-0.65cm]current page.north west) {%
        {\fontsize{5.3}{6.2}\selectfont\mediumfont\color{muted}#1 / NOTE DE PROJET}\par
        \vspace{0.18cm}
        {\fontsize{15.5}{16}\selectfont\displayfont\color{ink}#2}\par
        \vspace{0.20cm}\hairline\par\vspace{0.28cm}
        #3
        \vspace{0.08cm}\hairline\par\vspace{0.30cm}
        \RaggedRight{\fontsize{6.45}{8.0}\selectfont\lightfont\color{ink}#4}
      };
    \node[anchor=south west,inner sep=0,text=muted]
      at ([xshift=0.75cm,yshift=0.40cm]current page.south west)
      {\fontsize{4.8}{5.6}\selectfont\lightfont TEXTE + IMAGE / CONTEXTE ET INTENTION};
  \end{tikzpicture}\null
}

\newcommand{\galleryLabel}[2]{%
  \begin{tikzpicture}[remember picture,overlay]
    \fill[ink,opacity=0.90] ([xshift=0.45cm,yshift=0.42cm]current page.south west)
      rectangle ([xshift=7.4cm,yshift=1.45cm]current page.south west);
    \node[anchor=west,inner sep=0,text=white]
      at ([xshift=0.72cm,yshift=0.93cm]current page.south west)
      {\fontsize{5.2}{6.2}\selectfont\mediumfont #2\quad/\quad #1};
  \end{tikzpicture}%
}
\newcommand{\galleryFullPage}[3]{\clearpage\thispagestyle{empty}\fullBleedImage{#3}\galleryLabel{#1}{#2}\null}
\newcommand{\gallerySplitPage}[4]{%
  \clearpage\thispagestyle{empty}
  \begin{tikzpicture}[remember picture,overlay]
    \node[anchor=north west,inner sep=0] at (current page.north west){\includegraphics[width=12.7cm,height=14.8cm]{#3}};
    \node[anchor=north east,inner sep=0] at (current page.north east){\includegraphics[width=8.3cm,height=14.8cm]{#4}};
    \fill[acid] ([xshift=12.62cm]current page.south west) rectangle ([xshift=12.78cm]current page.north west);
  \end{tikzpicture}\galleryLabel{#1}{#2}\null
}
\newcommand{\galleryTriptychPage}[5]{%
  \clearpage\thispagestyle{empty}
  \begin{tikzpicture}[remember picture,overlay]
    \node[anchor=north west,inner sep=0] at (current page.north west){\includegraphics[width=12.9cm,height=14.8cm]{#3}};
    \node[anchor=north east,inner sep=0] at (current page.north east){\includegraphics[width=8.1cm,height=7.4cm]{#4}};
    \node[anchor=south east,inner sep=0] at (current page.south east){\includegraphics[width=8.1cm,height=7.4cm]{#5}};
    \fill[acid] ([xshift=12.82cm]current page.south west) rectangle ([xshift=12.98cm]current page.north west);
    \fill[acid] ([xshift=12.9cm,yshift=7.32cm]current page.south west) rectangle ([yshift=7.48cm]current page.south east);
  \end{tikzpicture}\galleryLabel{#1}{#2}\null
}
\newcommand{\galleryMosaicPage}[6]{%
  \clearpage\thispagestyle{empty}
  \begin{tikzpicture}[remember picture,overlay]
    \node[anchor=north west,inner sep=0] at (current page.north west){\includegraphics[width=12cm,height=14.8cm]{#3}};
    \node[anchor=north east,inner sep=0] at (current page.north east){\includegraphics[width=9cm,height=4.934cm]{#4}};
    \node[anchor=east,inner sep=0] at (current page.east){\includegraphics[width=9cm,height=4.934cm]{#5}};
    \node[anchor=south east,inner sep=0] at (current page.south east){\includegraphics[width=9cm,height=4.934cm]{#6}};
    \fill[acid] ([xshift=11.92cm]current page.south west) rectangle ([xshift=12.08cm]current page.north west);
  \end{tikzpicture}\galleryLabel{#1}{#2}\null
}

\begin{document}
\begin{titlepage}
\thispagestyle{empty}
\fullBleedImage{/home/alexandre/projets/atelier-archi/portfolio-pdf/output/crops/av-marne-maquette-cover-fee5285abae1.jpg}
\begin{tikzpicture}[remember picture,overlay]
  \fill[ink,opacity=0.36] (current page.south west) rectangle (current page.north east);
  \fill[ink,opacity=0.95] (current page.south west) rectangle ([xshift=7.6cm]current page.north west);
  \fill[acid] ([xshift=7.6cm]current page.south west) rectangle ([xshift=7.82cm]current page.north west);
  \node[anchor=north west,inner sep=0,text width=6cm]
    at ([xshift=0.75cm,yshift=-0.72cm]current page.north west) {%
      {\fontsize{5.2}{6.2}\selectfont\mediumfont\color{acid}ARCHITECTURE / DESIGN / ART}\par
      \vspace{1.05cm}
      {\fontsize{24}{22.5}\selectfont\displayfont\color{white}ALEXANDRE\\MATHIEU}\par
      \vspace{0.38cm}
      {\fontsize{10}{11.5}\selectfont\lightfont\color{white!72}PORTFOLIO\\2018?2026}
    };
  \node[anchor=south west,inner sep=0,text width=5.9cm]
    at ([xshift=0.75cm,yshift=0.72cm]current page.south west) {%
      {\fontsize{5.4}{6.6}\selectfont\lightfont\color{white!75}
      alexandre.mat+w@protonmail.com\par +33 6 58 21 53 00\par alexandre-mathieu-arch.github.io/works/}
    };
  \node[anchor=south east,rotate=90,inner sep=0,text=white]
    at ([xshift=-0.34cm,yshift=0.58cm]current page.south east)
    {\fontsize{4.8}{5.6}\selectfont\mediumfont SELECTED WORKS / 13 PROJECTS};
\end{tikzpicture}\null
\end{titlepage}

\clearpage\thispagestyle{empty}
\begin{tikzpicture}[remember picture,overlay]
  \fill[ink] (current page.south west) rectangle (current page.north east);
  \fill[acid] (current page.north west) rectangle ([xshift=0.22cm]current page.south west);
  \node[anchor=north west,inner sep=0,text width=18.5cm]
    at ([xshift=0.72cm,yshift=-0.65cm]current page.north west) {%
      {\fontsize{5.4}{6.3}\selectfont\mediumfont\color{acid}INDEX / SELECTED WORKS / 2018?2026}\par
      \vspace{0.26cm}
      {\fontsize{24}{23}\selectfont\displayfont\color{white}SOMMAIRE}\par
      \vspace{0.28cm}\textcolor{white!20}{\rule{\linewidth}{0.25pt}}
    };
  \node[anchor=south east,inner sep=0,text=white!10]
    at ([xshift=-0.28cm,yshift=0.1cm]current page.south east)
    {\fontsize{57}{54}\selectfont\displayfont 00};
\end{tikzpicture}
\vspace*{3.28cm}
\begin{multicols}{2}
\portfolioTocEntry{01}{La Marne}{2021}
\portfolioTocEntry{02}{Villa CN}{2024}
\portfolioTocEntry{03}{Base Vie}{2024}
\portfolioTocEntry{04}{Coliving Etel}{2023}
\portfolioTocEntry{05}{Porte-portable}{2023}
\portfolioTocEntry{06}{Atrium Organique}{2022}
\portfolioTocEntry{07}{Cap-Horn - Brule Architectes}{2022}
\portfolioTocEntry{08}{Maison Bretonne}{2020}
\portfolioTocEntry{09}{Portaledge}{2020}
\portfolioTocEntry{10}{Bodelio}{2019}
\portfolioTocEntry{11}{23 Plots}{2018}
\portfolioTocEntry{12}{230 Volts}{2018}
\portfolioTocEntry{13}{Les 3 Frères}{2018}
\end{multicols}

\hypertarget{project-01}{}

\projectIntroPage{01}{La Marne}{2021}{Lorient}{/home/alexandre/projets/atelier-archi/portfolio-pdf/output/crops/av-marne-maquette-intro-01-acaffaa86bdf.jpg}

\projectStoryPage{01}{La Marne}{\projectMetaItem{Typologie}{Résidentiel}
\projectMetaItem{Lieu}{Lorient}
\projectMetaItem{Pays}{France, Bretagne}
\projectMetaItem{Surface}{286 m²}
\projectMetaItem{Collaboration}{Pascal Debard (KLOUDE) \& Achraf Benbouzid (ABA)}
\projectMetaItem{Année}{2021}}{\leadText{Situé dans un jardin en intérieur d'îlot urbain, ce local commun a été transformé en résidence de 5 logements, du T2 au T4. Une renaissance architecturale privilégiant la lumière et l'intimité.}

Né de la transformation d’un ancien local commun en cœur d’îlot, La Marne propose une densification douce, à distance du tumulte urbain. Le projet conserve la qualité du jardin tout en redonnant une vraie valeur d’usage au bâti existant.\par La résidence rassemble cinq logements, du T2 au T4, pour une surface habitable totale de 286 m². L’intervention cherche un équilibre entre optimisation du foncier et qualité de vie, sans effacer le caractère calme du site.\par Les logements sont dessinés comme des refuges lumineux. Les rez-de-chaussée s’ouvrent sur des jardins privatifs, tandis que les logements à l’étage profitent de balcons et de vues dégagées vers le cœur d’îlot.\par Les stationnements sont intégrés avec discrétion afin de préserver la fluidité des cheminements. Chaque logement dispose d’un espace extérieur, d’une orientation travaillée et d’un rapport mesuré aux voisins.}{/home/alexandre/projets/atelier-archi/portfolio-pdf/output/crops/av-marne-rendu-ext-story-01-f862553ea86c.jpg}

\galleryMosaicPage{La Marne}{01 / 01}{/home/alexandre/projets/atelier-archi/portfolio-pdf/output/crops/existant-gallery-main-0-612a8dabe622.jpg}{/home/alexandre/projets/atelier-archi/portfolio-pdf/output/crops/existant-2-gallery-strip-0-1-65b77a914ae4.jpg}{/home/alexandre/projets/atelier-archi/portfolio-pdf/output/crops/existant-3-gallery-strip-0-2-4192d06cf6a8.jpg}{/home/alexandre/projets/atelier-archi/portfolio-pdf/output/crops/plan-masse-existant-gallery-strip-0-3-056c36ebe36c.jpg}


\galleryMosaicPage{La Marne}{01 / 02}{/home/alexandre/projets/atelier-archi/portfolio-pdf/output/crops/plan-niveau-axo-gallery-main-4-ee8034ecc30a.jpg}{/home/alexandre/projets/atelier-archi/portfolio-pdf/output/crops/rdc-gallery-strip-4-1-032908ec2861.jpg}{/home/alexandre/projets/atelier-archi/portfolio-pdf/output/crops/niveau1-gallery-strip-4-2-8954cfbb981e.jpg}{/home/alexandre/projets/atelier-archi/portfolio-pdf/output/crops/facade-ouest-gallery-strip-4-3-2d57050571ad.jpg}


\galleryTriptychPage{La Marne}{01 / 03}{/home/alexandre/projets/atelier-archi/portfolio-pdf/output/crops/facade-sud-gallery-main-8-bb3f1571ae45.jpg}{/home/alexandre/projets/atelier-archi/portfolio-pdf/output/crops/coupe-gallery-top-8-e0f69ec17212.jpg}{/home/alexandre/projets/atelier-archi/portfolio-pdf/output/crops/abris-velo-gallery-bottom-8-f7fae43e1368.jpg}


\hypertarget{project-02}{}

\projectIntroPage{02}{Villa CN}{2024}{Tahanaout, Marrakech / Livré / Réalisé}{/home/alexandre/projets/atelier-archi/portfolio-pdf/output/crops/villa-cn.01-intro-02-a2ead68cf107.jpg}

\projectStoryPage{02}{Villa CN}{\projectMetaItem{Typologie}{Résidentiel}
\projectMetaItem{Lieu}{Tahanaout, Marrakech}
\projectMetaItem{Pays}{Maroc}
\projectMetaItem{Surface}{650 m²}
\projectMetaItem{Phase}{Livré}
\projectMetaItem{Année}{2024}}{\leadText{Une villa familiale nichée dans les plaines bordant l’Atlas nord, entre murs épais, patios calmes et lumière maîtrisée.}

La Villa CN compose avec la chaleur, la lumière et l’épaisseur. Son architecture minérale installe des volumes protecteurs, ouverts par séquences vers le paysage et les patios.\par Les murs épais en terre et chaux donnent au projet son inertie et sa présence. Ils captent la chaleur du jour, tempèrent les espaces intérieurs et prolongent une expression constructive locale.\par Les ouvertures sont cadrées pour éviter l’exposition directe excessive et créer des ambiances contrastées. Les patios deviennent des pièces de calme, à la fois intérieures et extérieures.\par La villa recherche une forme de luxe silencieux : espaces généreux, matières naturelles, fraîcheur, ombre et vues choisies.}{/home/alexandre/projets/atelier-archi/portfolio-pdf/output/crops/villa-cn-atlas-story-02-7840db56884f.jpg}

\galleryFullPage{Villa CN}{02 / 01}{/home/alexandre/projets/atelier-archi/portfolio-pdf/output/crops/villa-cn02-gallery-full-0-1cbe95ffa537.jpg}


\hypertarget{project-03}{}

\projectIntroPage{03}{Base Vie}{2024}{}{/home/alexandre/projets/atelier-archi/portfolio-pdf/output/crops/base-vie-intro-03-7f947982b0c9.jpg}

\projectStoryPage{03}{Base Vie}{\projectMetaItem{Typologie}{Éphémère}
\projectMetaItem{Pays}{France}
\projectMetaItem{Année}{2024}}{\leadText{Une base vie de chantier pensée comme un lieu clair, robuste et confortable, au service des équipes sur site.}

Base Vie explore l’installation temporaire de chantier comme un espace d’accueil et de travail à part entière. Le projet cherche une organisation simple, lisible et rapide à mettre en œuvre.\par Les modules regroupent les fonctions essentielles du quotidien : réunion, pause, vestiaire et logistique. L’objectif est d’offrir un lieu efficace sans perdre la qualité d’ambiance nécessaire aux longues journées de chantier.\par La proposition privilégie des éléments robustes, démontables et faciles à entretenir. L’expression reste sobre pour laisser la priorité aux usages et à la durée de vie du dispositif.}{/home/alexandre/projets/atelier-archi/portfolio-pdf/output/crops/base-vie-story-03-606fa2b7d747.jpg}

\hypertarget{project-04}{}

\projectIntroPage{04}{Coliving Etel}{2023}{Etel / En cours / Projet}{/home/alexandre/projets/atelier-archi/portfolio-pdf/output/crops/coliving-etel-01-extension-intro-04-2119ff0daa4d.jpg}

\projectStoryPage{04}{Coliving Etel}{\projectMetaItem{Typologie}{Résidentiel}
\projectMetaItem{Lieu}{Etel}
\projectMetaItem{Pays}{France, Bretagne}
\projectMetaItem{Phase}{En cours}
\projectMetaItem{Année}{2023}}{\leadText{Une maison partagée à Étel, pensée pour associer intimité, espaces communs et relation apaisée au paysage côtier.}

À Étel, le projet s’appuie sur une maison existante pour imaginer un lieu de vie partagé, simple et lumineux. L’intervention cherche à préserver l’échelle domestique tout en clarifiant les usages collectifs.\par Les espaces communs sont placés au cœur de la maison, comme des lieux de rencontre quotidiens. Les chambres conservent une lecture plus intime, avec des seuils nets et des vues choisies.\par L’extension prolonge la maison sans la dominer. Elle apporte de la lumière, des surfaces utiles et une relation plus directe au jardin.}{/home/alexandre/projets/atelier-archi/portfolio-pdf/output/crops/coliving-etel-02-entree-story-04-6d7dfb840880.jpg}

\galleryMosaicPage{Coliving Etel}{04 / 01}{/home/alexandre/projets/atelier-archi/portfolio-pdf/output/crops/coliving-etel-03-situation-gallery-main-0-bfe558fe543d.jpg}{/home/alexandre/projets/atelier-archi/portfolio-pdf/output/crops/coliving-etel-04-masse-gallery-strip-0-1-6c13c623b413.jpg}{/home/alexandre/projets/atelier-archi/portfolio-pdf/output/crops/coliving-etel-05-rdc-gallery-strip-0-2-1eea635b77aa.jpg}{/home/alexandre/projets/atelier-archi/portfolio-pdf/output/crops/coliving-etel-06-etage-gallery-strip-0-3-dc7035c2ecb6.jpg}


\gallerySplitPage{Coliving Etel}{04 / 02}{/home/alexandre/projets/atelier-archi/portfolio-pdf/output/crops/coliving-etel-07-comble-gallery-left-4-81514ad4c243.jpg}{/home/alexandre/projets/atelier-archi/portfolio-pdf/output/crops/coliving-etel-08-chambre-haute-gallery-right-4-e22e30a56502.jpg}


\hypertarget{project-05}{}

\projectIntroPage{05}{Porte-portable}{2023}{Prototype / Réalisé}{/home/alexandre/projets/atelier-archi/portfolio-pdf/output/crops/porte-portable-okume-intro-05-ad845805432a.jpg}

\projectStoryPage{05}{Porte-portable}{\projectMetaItem{Typologie}{Design, Objet}
\projectMetaItem{Pays}{France}
\projectMetaItem{Phase}{Prototype}
\projectMetaItem{Année}{2023}}{\leadText{Série de prototypes pour smartphone en multiplis d'okoumé, cuir et inox. Ces objets ont été conçus pour accompagner la création d'un bureau en bois et acier, offrant une simplicité absolue sans mécanisme.}

Le porte-portable accompagne la création d’un bureau en bois et acier. Il prolonge la même recherche : un objet simple, précis, sans mécanisme superflu.\par Le multiplis d’okoumé forme la structure, le cuir protège le contact avec le téléphone et l’inox rappelle les piétements du bureau. Les assemblages restent lisibles et volontairement sobres.\par La géométrie permet quatre positions : deux inclinaisons et deux orientations, portrait ou paysage. L’objet est fabriqué à la main, avec une économie de moyens au service de la fonctionnalité.}{/home/alexandre/projets/atelier-archi/portfolio-pdf/output/crops/okume-2-story-05-d70157eac0ff.jpg}

\galleryFullPage{Porte-portable}{05 / 01}{/home/alexandre/projets/atelier-archi/portfolio-pdf/output/crops/okume-3-gallery-full-0-7cfb4ea484f7.jpg}


\hypertarget{project-06}{}

\projectIntroPage{06}{Atrium Organique}{2022}{Vannes / Concours privé (abandonné) / Abandonné}{/home/alexandre/projets/atelier-archi/portfolio-pdf/output/crops/night-shift-intro-06-08deb62b270a.jpg}

\projectStoryPage{06}{Atrium Organique}{\projectMetaItem{Typologie}{Tertiaire, Bureau}
\projectMetaItem{Lieu}{Vannes}
\projectMetaItem{Pays}{France, Bretagne}
\projectMetaItem{Phase}{Concours privé (abandonné)}
\projectMetaItem{Année}{2022}}{\leadText{Des bulles en toile tendue introduisent des espaces calmes et confidentiels au cœur d’un atrium bancaire très orthogonal.}

Atrium Organique intervient dans un grand hall bancaire marqué par une trame orthogonale très présente. Le projet introduit des volumes plus souples, capables d’apporter du calme et de nouvelles situations d’usage sans alourdir l’espace existant.\par Des structures métalliques légères, habillées de toiles translucides, dessinent des bulles programmatiques qui semblent flotter dans l’atrium. Chacune accueille une fonction précise : réunion, détente ou échange informel.\par La toile filtre la lumière et laisse deviner les silhouettes, tout en offrant un degré d’intimité. La géométrie organique adoucit la lecture du hall et transforme un grand volume résonant en espace plus feutré.\par Les formes sont générées avec Rhinocéros 3D et Grasshopper, puis intégrées au modèle BIM via Bridge Archicad. Les images de travail sont produites avec Twinmotion et Photoshop.}{/home/alexandre/projets/atelier-archi/portfolio-pdf/output/crops/night-shift-bis-story-06-cfe14fdfe1dc.jpg}

\galleryFullPage{Atrium Organique}{06 / 01}{/home/alexandre/projets/atelier-archi/portfolio-pdf/output/crops/perspective-gallery-full-0-a581fb3f128b.jpg}


\hypertarget{project-07}{}

\projectIntroPage{07}{Cap-Horn - Brule Architectes}{2022}{Quimper / Réalisé}{/home/alexandre/projets/atelier-archi/portfolio-pdf/output/crops/cap-horn-01-perspective-intro-07-cff8c2934619.jpg}

\projectStoryPage{07}{Cap-Horn - Brule Architectes}{\projectMetaItem{Typologie}{Tertiaire, Bureau}
\projectMetaItem{Lieu}{Quimper}
\projectMetaItem{Pays}{France, Bretagne}
\projectMetaItem{Surface}{850 m²}
\projectMetaItem{Phase}{Réalisé}
\projectMetaItem{Année}{2022}}{Cap-Horn est un projet d’agence développé au sein de Brulé Architectes Associés. Situé sur une parcelle étroite à Quimper, il répond par une organisation verticale claire, capable de superposer plusieurs usages sans brouiller la lecture du bâtiment.\par Le projet rassemble un socle de stationnement, des plateaux destinés aux professions libérales et les espaces de l’agence dans les derniers niveaux. Cette stratification donne au bâtiment une logique simple : les usages les plus ouverts restent au contact de la ville, tandis que les espaces de travail prennent de la hauteur.\par Les derniers niveaux offrent une position plus calme, à distance de la rue, avec des vues dégagées sur l’horizon urbain. Le toit-terrasse prolonge cette qualité d’usage par une salle de réunion panoramique.\par Les circulations verticales sont externalisées afin de libérer les plateaux et de préserver leur souplesse d’aménagement. L’ascenseur et l’escalier deviennent ainsi des éléments structurants de la façade autant que des outils d’organisation intérieure.}{/home/alexandre/projets/atelier-archi/portfolio-pdf/output/crops/cap-horn-02-terrasse-story-07-d657747f658d.jpg}

\galleryMosaicPage{Cap-Horn - Brule Architectes}{07 / 01}{/home/alexandre/projets/atelier-archi/portfolio-pdf/output/crops/cap-horn-03-elevation-ouest-gallery-main-0-2b234839d934.jpg}{/home/alexandre/projets/atelier-archi/portfolio-pdf/output/crops/cap-horn-04-principe-facade-gallery-strip-0-1-96ba8213c5f8.jpg}{/home/alexandre/projets/atelier-archi/portfolio-pdf/output/crops/cap-horn-05-detail-coupe-gallery-strip-0-2-a3d76db24e7b.jpg}{/home/alexandre/projets/atelier-archi/portfolio-pdf/output/crops/cap-horn-06-plan-masse-gallery-strip-0-3-a35847ddc9fa.jpg}


\hypertarget{project-08}{}

\projectIntroPage{08}{Maison Bretonne}{2020}{La Trinité-sur-Mer / Livré / Réalisé}{/home/alexandre/projets/atelier-archi/portfolio-pdf/output/crops/maison-bretonne-ext-intro-08-8bf67cc58aa9.jpg}

\projectStoryPage{08}{Maison Bretonne}{\projectMetaItem{Typologie}{Résidentiel}
\projectMetaItem{Lieu}{La Trinité-sur-Mer}
\projectMetaItem{Pays}{France, Bretagne}
\projectMetaItem{Surface}{150 m²}
\projectMetaItem{Phase}{Livré}
\projectMetaItem{Collaboration}{Pascal Debard (KLOUDE)}
\projectMetaItem{Année}{2020}}{\leadText{Dialogue entre héritage néo-breton et légèreté contemporaine : rénovation et extension d'une maison côtière en structure bois et zinc blanc.}

Située près du Golfe du Morbihan, la Maison Bretonne engage un dialogue entre héritage néo-breton et écriture contemporaine. Le projet conserve la présence du volume existant tout en l’ouvrant à de nouveaux usages.\par La rénovation complète s’accompagne d’une extension en structure bois et zinc blanc. L’enduit clair unifie les volumes et prolonge l’image sobre des maisons côtières.\par L’isolation thermique par l’extérieur améliore la performance énergétique et donne une lecture plus calme de l’ensemble. Les espaces intérieurs gagnent en lumière, en continuité et en rapport au paysage.\par Réalisé en indépendant, ce projet installe une attention durable au contexte, à la matière et à la précision constructive.}{/home/alexandre/projets/atelier-archi/portfolio-pdf/output/crops/maison-bretonne-int-story-08-5e80077d955f.jpg}

\galleryMosaicPage{Maison Bretonne}{08 / 01}{/home/alexandre/projets/atelier-archi/portfolio-pdf/output/crops/maison-bretonne-situation-gallery-main-0-5bc820d704fd.jpg}{/home/alexandre/projets/atelier-archi/portfolio-pdf/output/crops/maison-bretonne-rdc-gallery-strip-0-1-b5dba220ddf2.jpg}{/home/alexandre/projets/atelier-archi/portfolio-pdf/output/crops/maison-bretonne-etage-gallery-strip-0-2-802166384a71.jpg}{/home/alexandre/projets/atelier-archi/portfolio-pdf/output/crops/maison-bretonne-facade-sud-gallery-strip-0-3-13a82f8d54da.jpg}


\galleryMosaicPage{Maison Bretonne}{08 / 02}{/home/alexandre/projets/atelier-archi/portfolio-pdf/output/crops/maison-bretonne-facade-ouest-gallery-main-4-72114268ba04.jpg}{/home/alexandre/projets/atelier-archi/portfolio-pdf/output/crops/maison-bretonne-extrait-coupes-gallery-strip-4-1-413799435b78.jpg}{/home/alexandre/projets/atelier-archi/portfolio-pdf/output/crops/maison-bretonne-extrait-coupe-gallery-strip-4-2-a0b483c1bc81.jpg}{/home/alexandre/projets/atelier-archi/portfolio-pdf/output/crops/maison-bretonne-avant-gallery-strip-4-3-300e3e0e56e7.jpg}


\galleryTriptychPage{Maison Bretonne}{08 / 03}{/home/alexandre/projets/atelier-archi/portfolio-pdf/output/crops/maison-bretonne-photogrammetrie-gallery-main-8-b1d59e420a1f.jpg}{/home/alexandre/projets/atelier-archi/portfolio-pdf/output/crops/non-visible-detail-volet-roulant-gallery-top-8-e3ba456a9366.jpg}{/home/alexandre/projets/atelier-archi/portfolio-pdf/output/crops/non-visible-facades-gallery-bottom-8-569f2a365572.jpg}


\hypertarget{project-09}{}

\projectIntroPage{09}{Portaledge}{2020}{Lorient / Études / En cours}{/home/alexandre/projets/atelier-archi/portfolio-pdf/output/crops/studiotetu-intro-09-6a83a0d127dc.jpg}

\projectStoryPage{09}{Portaledge}{\projectMetaItem{Typologie}{Résidentiel}
\projectMetaItem{Lieu}{Lorient}
\projectMetaItem{Pays}{France, Bretagne}
\projectMetaItem{Surface}{25 m²}
\projectMetaItem{Phase}{Études}
\projectMetaItem{Année}{2020}}{\leadText{Un studio suspendu compact, pensé comme un refuge de travail minimal, entre verticalité, silence et concentration.}

Portaledge imagine un volume de travail réduit à l’essentiel. Le projet s’intéresse à la sensation d’abri, à la maîtrise de l’encombrement et à la qualité d’un espace isolé.\par Le studio privilégie une organisation compacte : un poste de travail, quelques rangements, une assise et une ouverture généreuse. Chaque élément participe à la clarté du lieu.\par La suspension donne au projet son caractère. Elle libère le sol, renforce la relation au paysage et transforme le studio en observatoire calme.}{/home/alexandre/projets/atelier-archi/portfolio-pdf/output/crops/studiotetu-story-09-6368cd5cfd1e.jpg}

\hypertarget{project-10}{}

\projectIntroPage{10}{Bodelio}{2019}{Lorient / Concours / Concours perdu}{/home/alexandre/projets/atelier-archi/portfolio-pdf/output/crops/bodelio-pers-intro-10-40a269effcb6.jpg}

\projectStoryPage{10}{Bodelio}{\projectMetaItem{Typologie}{Résidentiel}
\projectMetaItem{Lieu}{Lorient}
\projectMetaItem{Pays}{France, Bretagne}
\projectMetaItem{Surface}{4200 m²}
\projectMetaItem{Phase}{Concours}
\projectMetaItem{Collaboration}{DDL Architectes}
\projectMetaItem{Année}{2019}}{\leadText{Un projet de logements collectifs inscrit dans la reconversion du site hospitalier Bodelio à Lorient.}

Situé sur l’ancien site de l’hôpital Bodelio, le projet accompagne la transformation d’un morceau de ville en quartier habité. L’enjeu est de densifier le centre de Lorient tout en préservant des vues, des respirations et une vraie qualité d’usage.\par Les volumes sont fragmentés pour retrouver l’échelle des îlots voisins. Cette organisation multiplie les orientations, ménage des percées visuelles et évite l’effet de masse d’un front bâti unique.\par Chaque logement bénéficie d’un prolongement extérieur généreux, pensé comme une pièce supplémentaire. Les jardins et espaces communs au rez-de-chaussée installent une relation simple entre intimité domestique et vie collective.\par Le projet privilégie une matérialité durable : plateaux en CLT, attiques en ossature bois, balcons portés par une structure apparente. Le stationnement souterrain libère le cœur d’îlot pour le paysage.}{/home/alexandre/projets/atelier-archi/portfolio-pdf/output/crops/maquette-story-10-161742d541e2.jpg}

\hypertarget{project-11}{}

\projectIntroPage{11}{23 Plots}{2018}{Caen / Concours}{/home/alexandre/projets/atelier-archi/portfolio-pdf/output/crops/plot-logement-caen-intro-11-4b1d7f4dac5b.jpg}

\projectStoryPage{11}{23 Plots}{\projectMetaItem{Typologie}{Résidentiel}
\projectMetaItem{Lieu}{Caen}
\projectMetaItem{Pays}{France, Normandie}
\projectMetaItem{Phase}{Concours}
\projectMetaItem{Collaboration}{DDL architectes Lorient}
\projectMetaItem{Année}{2018}}{\leadText{Concours pour un ensemble de 23 plots de logements à Caen, organisé autour d’une densité ouverte et paysagère.}

Le projet 23 Plots interroge la densité résidentielle à grande échelle. Plutôt qu’un front continu, l’opération répartit les logements en volumes autonomes pour ménager des vues, des passages et des respirations.\par Les plots s’inscrivent dans une trame paysagère qui structure les cheminements et les espaces communs. Cette organisation donne une lecture plus douce de l’ensemble et améliore les relations entre logements.\par Chaque bâtiment cherche une orientation favorable et un rapport direct aux extérieurs. La densité devient un support de qualité d’usage plutôt qu’un simple objectif quantitatif.}{/home/alexandre/projets/atelier-archi/portfolio-pdf/output/crops/plot-logement-caen-story-11-ffecab57609e.jpg}

\hypertarget{project-12}{}

\projectIntroPage{12}{230 Volts}{2018}{Prototype / Réalisé}{/home/alexandre/projets/atelier-archi/portfolio-pdf/output/crops/montage-prise-intro-12-378a03dd8a85.jpg}

\projectStoryPage{12}{230 Volts}{\projectMetaItem{Typologie}{Design}
\projectMetaItem{Pays}{France}
\projectMetaItem{Phase}{Prototype}
\projectMetaItem{Collaboration}{École d'Architecture de Nantes}
\projectMetaItem{Année}{2018}}{\leadText{Un dispositif d’alimentation mobile et discret, conçu pour libérer les usages dans les studios de l’École d’Architecture de Nantes.}

Le projet 230 Volts transforme une contrainte technique en objet mobile. Dans les studios de l’ENSA Nantes, les prises fixées aux chemins de câbles limitent les déplacements et imposent une position de travail.\par Le dispositif associe une alimentation descendante, un rail de translation et un panier vide-poche. Il rapproche l’énergie du poste de travail sans encombrer le sol et sécurise les appareils légers pendant la charge.\par L’objet reprend la logique constructive du bâtiment de Lacaton \& Vassal : matériaux simples, assemblage lisible, présence mesurée. La tôle perforée, le métal et les détails de suspension prolongent l’atmosphère industrielle des ateliers.\par La prise se déplace horizontalement et verticalement selon le besoin. Elle accompagne les gestes quotidiens avec une économie de moyens volontaire.}{/home/alexandre/projets/atelier-archi/portfolio-pdf/output/crops/realisation-story-12-670ea59e0f48.jpg}

\galleryMosaicPage{230 Volts}{12 / 01}{/home/alexandre/projets/atelier-archi/portfolio-pdf/output/crops/prise-multi-gallery-main-0-483c6e9ff64e.jpg}{/home/alexandre/projets/atelier-archi/portfolio-pdf/output/crops/axonometrie-gallery-strip-0-1-5b8159cf2e3e.jpg}{/home/alexandre/projets/atelier-archi/portfolio-pdf/output/crops/maquette-gallery-strip-0-2-8d522ffdaa49.jpg}{/home/alexandre/projets/atelier-archi/portfolio-pdf/output/crops/maquette-papier-gallery-strip-0-3-8da8d9c6621b.jpg}


\galleryMosaicPage{230 Volts}{12 / 02}{/home/alexandre/projets/atelier-archi/portfolio-pdf/output/crops/vue-elevation-gallery-main-4-4d199941c10b.jpg}{/home/alexandre/projets/atelier-archi/portfolio-pdf/output/crops/coupe-projet-230-gallery-strip-4-1-bfb0e96c8f3a.jpg}{/home/alexandre/projets/atelier-archi/portfolio-pdf/output/crops/schema-montage-gallery-strip-4-2-607ea6f5e853.jpg}{/home/alexandre/projets/atelier-archi/portfolio-pdf/output/crops/materialite-oc-gallery-strip-4-3-365f352d3062.jpg}


\galleryFullPage{230 Volts}{12 / 03}{/home/alexandre/projets/atelier-archi/portfolio-pdf/output/crops/type_onze_tole_perforee-gallery-full-8-f45138f69e27.jpg}


\hypertarget{project-13}{}

\projectIntroPage{13}{Les 3 Frères}{2018}{Saint-Brieuc}{/home/alexandre/projets/atelier-archi/portfolio-pdf/output/crops/3frereslg-lota-pers-intro-13-869cab90b229.jpg}

\projectStoryPage{13}{Les 3 Frères}{\projectMetaItem{Typologie}{Résidentiel}
\projectMetaItem{Lieu}{Saint-Brieuc}
\projectMetaItem{Pays}{France, Bretagne}
\projectMetaItem{Surface}{565 m²}
\projectMetaItem{Collaboration}{Agence DDL Lorient}
\projectMetaItem{Année}{2018}}{\leadText{Sept logements du T2 au T4 à Saint-Brieuc, avec stationnement privatif en sous-sol et jardins en cœur d’îlot.}

Le projet Les 3 Frères s’inscrit dans un tissu résidentiel de Saint-Brieuc. Il propose une opération compacte, organisée autour de logements traversants ou bien orientés, avec une attention portée aux prolongements extérieurs.\par Le bâtiment compose avec l’échelle de la rue et le calme de l’intérieur d’îlot. Les stationnements sont placés en sous-sol pour libérer les espaces extérieurs et préserver la qualité des jardins.\par La matérialité associe béton, verre et tôle perforée. Les garde-corps et filtres visuels donnent de l’intimité aux logements tout en maintenant une façade claire et régulière.}{/home/alexandre/projets/atelier-archi/portfolio-pdf/output/crops/3frereslg-lota-pers-jardins-story-13-12d09c694d2f.jpg}

\gallerySplitPage{Les 3 Frères}{13 / 01}{/home/alexandre/projets/atelier-archi/portfolio-pdf/output/crops/3frereslg-lotb-pers-gallery-left-0-25c5b3a394c6.jpg}{/home/alexandre/projets/atelier-archi/portfolio-pdf/output/crops/3frereslg-planmasse-gallery-right-0-db7212e379ca.jpg}

\end{document}
